\section{Determining the covariance below \texorpdfstring{\(c=1\)}{c=1}}\label{sec:normalization}

Throughout this section \(0<c<1\), so \(2\pi/3<\lambda<\pi\).
Our objective is to determine the orbit mass \(\kappa\) defined in
Section~\ref{sec:angular}. Two different finite systems give opposite
bounds. A homogeneous chain at fixed particle density gives a lower
bound through its response to a seam twist. A dense array of real
inhomogeneities gives an upper bound through the change of its vacuum
vector under the same twist. Neither comparison uses positive association.

\begin{proposition}[The covariance constant below one]\label{prop:normalization}
With the row and height conventions of Sections~\ref{sec:finite}
and~\ref{sec:angular},
\begin{equation}\label{eq:kappa-value}
 \kappa=\frac1{2\pi(\pi-\lambda)}.
\end{equation}
Consequently a limiting covariance \(-2\kappa\log|x-y|\) has squared
multiplier \(4\pi\kappa=2/(\pi-\lambda)=\sigma(c)^2\) relative to
\eqref{eq:green}.
\end{proposition}

\subsection{Homogeneous stiffness at fixed density}

For \(Q<\infty\), let \(K_Q\) denote convolution with \(K\) restricted to
\([-Q,Q]\), and write
\[
 r_Q=(I+K_Q)^{-1}s,\qquad H_Q=(I+K_Q)^{-1}\bbone.
\]
These are the density and dressed constant of Lemma~\ref{lem:root-limits}.
Choose the homogeneous finite chain with particle number
\(n/L\to\int_{-Q}^{Q}r_Q(u)\,\dd u\). The first limit in this subsection
is \(L\to\infty\) at this fixed \(Q\); only afterwards does \(Q\to\infty\).
After both limits, the equal-time local state is the balanced plane
state, by Proposition~\ref{prop:word-kernel} and
Proposition~\ref{prop:physical-plane}. At fixed \(Q\), the filling
remains strictly below one half.

Insert the auxiliary seam \(e^{itZ}\), and let \(\mathcal H_L(t)\) be the
endpoint Hamiltonian before subtracting its ground energy. Denote its
continued sector ground energy by \(e_L(t)\), its unit ground vector at
zero twist by \(\Omega_L\), and
\[
 E_{\mathrm{exc},L}
   =\mathcal H_L(0)-e_L(0)
\]
on that magnetization sector. The untwisted sector minimum is simple.
All inverse-energy expressions below are restricted to
\(\Omega_L^\perp\) in that sector.

At the swap endpoint, the auxiliary spin is carried from site to site.
The seam therefore twists the closing link. Conjugating by site rotations
whose angles increase in increments \(t/L\) distributes this twist over
all links, since each link conserves the sum of its two spins. Let
\(b_j\) and \(d_j\) be respectively the first and second derivatives of
the local link under the single-link twist, with a common choice of
orientation. Put \(\mathcal J_L=\sum_j b_j\) and
\(D_{0,L}=L^{-1}\sum_j d_j\).
Ordinary second-order perturbation theory then gives
\begin{equation}\label{eq:kohn}
 L e_L''(0)=\langle D_{0,L}\rangle
 -\frac2L
 \left\|E_{\mathrm{exc},L}^{-1/2}
       (I-|\Omega_L\rangle\langle\Omega_L|)\mathcal J_L\Omega_L\right\|^2 .
\end{equation}
This boundary-twist curvature identity has a standard spin-chain
formulation in Shastry and Sutherland~\cite[pp.~243--244]{ShastrySutherland1990},
who trace the method to Kohn.
Indeed, after distributing the seam,
\(\mathcal H_L'(0)=L^{-1}\mathcal J_L\) and
\(\mathcal H_L''(0)=L^{-2}\sum_jd_j\).
The sign chosen for \(t\) has no effect on \eqref{eq:kohn}.
For either orientation, the second link derivative is
\(d_j=-[Z_j,[Z_j,h_{j,j+1}]]\). Thus the completed local-state
limit of \(\langle D_{0,L}\rangle\) is precisely the constant
\(D_0\) in Lemma~\ref{lem:sum-rule}.

\begin{lemma}[Homogeneous stiffness]\label{lem:stiffness}
The iterated homogeneous stiffness satisfies
\begin{equation}\label{eq:stiffness-value}
 \lim_{Q\to\infty}\lim_{L\to\infty}L e_L''(0)
      =\frac2{\pi-\lambda}.
\end{equation}
\end{lemma}
\begin{proof}
The root-dependent part of the endpoint energy is
\(\sum_j e(y_j)\), where
\[
 e(u)=-2\pi\eta s(u).
\]
Differentiating the counting equations at zero twist gives
\[
 G y'=\frac{\bbone}{\pi},\qquad
 (G y'')_j=-L s'(y_j)(y'_j)^2
      +\sum_kK'(y_j-y_k)(y'_j-y'_k)^2.
\]
All derivatives here are taken on the regular real-root branch at fixed
finite \(L\). Lemma~\ref{lem:gaudin-inverse} and the fixed-\(Q\) quadratures
imply, at the nodes,
\[
 L y_j'\longrightarrow W(y_j),\qquad
 W=\frac{H_Q}{\pi r_Q}.
\]
Set \(\varepsilon_Q=(I+K_Q)^{-1}e'\).
The symmetry of \(G\) gives
\(L G^{-1}e'\to\varepsilon_Q/r_Q\).
Substituting the equation for \(y''\) into the second energy derivative
and expanding its square therefore gives
\begin{align}
 \lim_{L\to\infty}L e_L''(0)
 &=\int_{-Q}^{Q}
   \left\{(r_Q\varepsilon_Q'-\varepsilon_Qr_Q')W^2
       +2\varepsilon_Q W(r_QW)'\right\}\,\dd u \notag\\
 &=2\bigl(\varepsilon_Qr_QW^2\bigr)(Q).
 \label{eq:stiffness-boundary}
\end{align}
Here is the cancellation in the first equality. The direct squared
displacement term is \(\int r_Q e''W^2\).
The term \(-s'+K_Q'r_Q\) in the implicit derivative is
\(-r_Q'\), so it contributes \(-\int\varepsilon_Qr_Q'W^2\).
Transposing the other \(K_Q'\) term replaces \(e''\) by
\(\varepsilon_Q'\). The remaining cross term uses
\((I+K_Q)(r_QW)=1/\pi\).
The integrand is now the derivative of
\(\varepsilon_Qr_QW^2\). The density and \(W\) are even and
\(\varepsilon_Q\) is odd, proving the boundary expression.

We next evaluate its boundary factors, rather than differentiating a
bulk asymptotic. Write
\[
 r(u)=\frac1{2\lambda\cosh(u/\eta)}.
\]
If \(g\) solves the full-line equation with the relevant source, its
restriction is corrected to the interval solution by
\[
 g\ \longmapsto\
 g+(I+K_Q)^{-1}K_{\mathrm{out}}g,
\]
where \(K_{\mathrm{out}}g\) integrates over \(\R\setminus[-Q,Q]\).
Apply this map to \(r\) and to \(-2\pi\eta r'\).
On the outer tail near \(Q\) their ratio tends uniformly to \(2\pi\).
The contribution of the opposite tail is negligible relative to \(r(Q)\):
insert a sufficiently small exponential weight in the convergent Neumann
series for \(K_Q\), and use the separation of the two ends.
The same comparison gives \(r_Q(Q)\asymp r(Q)\), hence
\[
 \frac{\varepsilon_Q(Q)}{r_Q(Q)}\longrightarrow2\pi.
\]

After translating the endpoint \(Q\) to zero, the Neumann series for
\(H_Q\) converges term by term to its half-line series and is dominated
by the geometric series with ratio \(\|K\|_1<1\). Thus \(H_Q(Q)\) has
a limit. Symmetry of the resolvent, or differentiation of its series,
gives the exact identity
\begin{equation}\label{eq:dressed-constant-boundary}
 \frac{\dd}{\dd Q}\int_{-Q}^{Q}H_Q(u)\,\dd u=2H_Q(Q)^2.
\end{equation}
The bulk series also gives
\[
 \frac1{2Q}\int_{-Q}^{Q}H_Q(u)\,\dd u
       \longrightarrow\frac1{1+\widehat K(0)}.
\]
Since the right side of \eqref{eq:dressed-constant-boundary} has a limit,
averaging that identity identifies
\[
 H_Q(Q)^2\longrightarrow\frac1{1+\widehat K(0)}
       =\frac{\pi}{2(\pi-\lambda)}.
\]
Finally \eqref{eq:stiffness-boundary} equals
\(2\varepsilon_Q(Q)H_Q(Q)^2/[\pi^2r_Q(Q)]\).
The two boundary limits yield \eqref{eq:stiffness-value}.
\end{proof}

\subsection{An infrared comparison with the plane}

The inverse in \eqref{eq:kohn} is not a local observable.
We therefore regularize it before passing to the plane.
For \(u>0\) put
\[
 g_u(E)=\frac{1-e^{-uE^2}}{E},\qquad g_u(0)=0.
\]
On nonnegative energies \(0\le g_u(E)\le E^{-1}\).
Extend \(g_u\) as an odd function to the real line when using its
Fourier representation.

\begin{lemma}[Infrared comparison]\label{lem:infrared}
With \(\kappa\) as in Section~\ref{sec:angular},
\begin{equation}\label{eq:kappa-lower}
 \kappa\ge\frac1{2\pi(\pi-\lambda)}.
\end{equation}
\end{lemma}
\begin{proof}
Replace the inverse-energy quadratic form in \eqref{eq:kohn} by
the smaller form with \(g_u(E_{\mathrm{exc},L})\), keeping \(u\) fixed.
Translation invariance rewrites its factor \(2/L\) as the sum over
link separations of twice a symmetrized single-source pairing.

An explicit real-time representation is
\[
 g_u(E)=\int_0^\infty
       \operatorname{erfc}\!\left(\frac{t}{2\sqrt u}\right)
       \sin(tE)\,\dd t,\qquad
 \operatorname{erfc}(x)=\frac2{\sqrt\pi}\int_x^\infty e^{-s^2}\,\dd s.
\]
Integration by parts gives this identity. Its sine kernel is bounded at
zero, integrable, and has a Gaussian tail at large real times.
Oddness expresses the symmetrized pairing through expectations of
real-time commutators of the local currents.
The endpoint-Hamiltonian identification in
Lemma~\ref{lem:endpoint-Hamiltonian} makes this the same real-time evolution
as that generated by the finite-range link Hamiltonian.
For link separation \(r\), split the time integral at a fixed small
multiple of \(r\). In the shorter-time part the finite-range
Lieb--Robinson estimate gives exponential decay in \(r\); in the
longer-time part the filter tail is summable without using locality
\cite{LiebRobinson1972}. The estimates are uniform in the circle distance.
For bounded times the local evolutions are norm approximated on a bounded
neighborhood, so local convergence of the ground states gives convergence
of each such commutator.

Let \(\beta_b\) be the spectral measure of the plane vector
\(b_0\Omega\) for the commuting angular rows; the endpoint energy
\(E\) and shift phase \(p\) are functions on this spectrum. The preceding summability shows that
\[
 2\operatorname{sym}\bigl((p)_*(g_u(E)\beta_b)\bigr)
       =\mathcal S_u(p)\,\frac{\dd p}{2\pi}
\]
has a continuous density. The finite regularized term converges to
\(\mathcal S_u(0)\), first as \(L\to\infty\) and then as
\(Q\to\infty\). Throughout this passage \(u\) is fixed.
Consequently \eqref{eq:kohn} and Lemma~\ref{lem:stiffness} imply
\begin{equation}\label{eq:regularized-stiffness}
 \frac2{\pi-\lambda}\le D_0-\mathcal S_u(0).
\end{equation}
The current and local spin preserve the magnetization sector, so the
nonnegativity used in this comparison is sector nonnegativity. The
real-time commutator estimate itself needs no ground-state hypothesis.

At nonzero momentum, the continuity equation
\[
 i[E,Z_j]=b_j-b_{j-1}
\]
identifies the current spectral density. On the vacuum,
\(\|Z_0\Omega\|^2(\dd f)=|1-f(i)|^2\mu(\dd f)\).
The continuity equation implies, away from \(p=0\), that the
energy-weighted current measure satisfies
\(g_u(E)\beta_b=E(1-e^{-uE^2})\mu\).
Indeed the spin measure carries \(|1-f(i)|^2\), which cancels the
squared shift difference in the continuity equation. Subtract this
identity from the sum rule in Lemma~\ref{lem:sum-rule}. On the
punctured phase circle this gives
\begin{equation}\label{eq:filtered-density}
 \bigl(D_0-\mathcal S_u(p)\bigr)\frac{\dd p}{2\pi}
 =2\operatorname{sym}\bigl((p)_*(Ee^{-uE^2}\mu)\bigr).
\end{equation}
The equality extends to the limit \(p\to0\) because the left side
is continuous. Indeed, the summable commutator bounds show that the
fixed-\(u\) susceptibility measure has an absolutely convergent spatial
Fourier series, hence no atom at phase zero. The regulator annihilates
zero energy, and the \(f=1\) current component was removed in the
construction of \(\mu\). Thus the off-zero division has omitted no
additional contribution to this continuous density.

On a dilation orbit write \(p(s)=\arg f(ie^s)\).
Its absolute derivative is \(E\), and its winding number is \(n_f\)
when the Blaschke degree is \(2n_f\). For small nonzero \(p\) of either
sign, changing variables from \(s\) to \(p\) in
\eqref{eq:filtered-density} leaves \(n_f\) crossings, weighted by
\(e^{-uE^2}\), and the density convention contributes \(4\pi\).
As \(p\to0\), precisely one crossing escapes to the corresponding
orbit end, where \(E\to0\); the other crossings have positive limiting
energy. The sum is bounded by \(n_f\), and
\(\int n_f\,\dd\nu<\infty\) by the sum rule. Dominated convergence,
first at \(p\to0\) and then at \(u\to\infty\), therefore gives
\[
 \lim_{u\to\infty}\bigl(D_0-\mathcal S_u(0)\bigr)=4\pi\kappa.
\]
Combining this with \eqref{eq:regularized-stiffness} proves
\eqref{eq:kappa-lower}.
\end{proof}

\subsection{The exact overlap to be estimated}

We now use the other finite system: real site parameters on an array
of spacing \(d/2\), of length \(B=Ld/2\), at half filling.
Let \(\Omega_{L,t}\) be the unit real-root vector with seam twist
\(e^{itZ}\), continued near zero, and put
\begin{equation}\label{eq:fidelity-cost}
 I_L=\left\|(I-|\Omega_L\rangle\langle\Omega_L|)
           \partial_t\Omega_{L,t}\big|_{t=0}\right\|^2.
\end{equation}
These staggered vectors need no Perron interpretation.

Let \(A\) and \(C\) be the root lists at zero and at \(t\).
Use \([x]=\sinh x\), and let \(M_{AC}\) be the reduced scalar-product
matrix of Section~\ref{sec:finite}, with on-shell roots \(A\) and column
arguments \(C\). In \(M_{CA}\), use the twisted on-shell roots \(C\)
and their twisted Bethe ratio \(e^{-2it}q_C\). At \(t=0\) both coincident
matrices are the appropriate Gaudin matrices. Lemmas~\ref{lem:scalar-product}
and~\ref{lem:hilbert-dual}, followed by the hyperbolic Cauchy identity, give
\begin{equation}\label{eq:exact-overlap}
 |\langle\Omega_{L,0},\Omega_{L,t}\rangle|^2
 =
 \prod_{j,k}
 \frac{[A_j-C_k+i\lambda][C_j-A_k+i\lambda]}
      {[A_j-A_k+i\lambda][C_j-C_k+i\lambda]}
 \frac{\det M_{AC}\det M_{CA}}{\det G_A\det G_C}.
\end{equation}
This is the identity \eqref{eq:overlap-determinant-finite}: its
one-set factors cancel, including those converting the algebraic duals
into Hilbert duals. It therefore measures the overlap of unit vectors.

For each finite chain, the analytic implicit-function theorem applies
to the counting equations because \(G\) is nonsingular.
The roots remain distinct, the vector remains nonzero and its squared
normalized overlap is nonzero on a sufficiently small neighborhood of
zero. The local logarithm with value zero at \(t=0\) is therefore defined.
Differentiating a unit vector shows
\begin{equation}\label{eq:overlap-quadratic}
 \log|\langle\Omega_{L,0},\Omega_{L,t}\rangle|^2
       =-I_Lt^2+O_L(t^3).
\end{equation}
Only its quadratic coefficient is used. There is no need for a
complex neighborhood whose radius is uniform in \(L\).

\begin{lemma}[Staggered overlap asymptotic]\label{lem:overlap-asymptotic}
In the order \(B\to\infty\), then \(d\downarrow0\),
\begin{equation}\label{eq:overlap-density}
 \frac{I_L}{B}\longrightarrow\frac1{2\lambda(\pi-\lambda)}.
\end{equation}
The limit is understood in the double-limit convention
\eqref{eq:array-limit-convention}.
\end{lemma}

We first use this asymptotic to bound \(\kappa\). Its proof occupies
Sections~\ref{sec:finite-trace}--\ref{sec:overlap-symbol} below.

\subsection{Packing angular windows}

We connect the staggered susceptibility to the plane by detecting
separated angular ranges in the spectrum of its derivative vector.
Let \(\zeta_L\) be the spectral measure of the orthogonal derivative vector
in \eqref{eq:fidelity-cost}, for the simultaneous self-adjoint rows
\(P_L(e^s)\), divided by their designated untwisted vacuum eigenvalues.
Its total mass is \(I_L\). Denote its spectral coordinates by \(f_s\).
The differentiated eigen-equation expresses multiplication of that vector
by \(1-f_s\) as the marked row vector, with its vacuum component removed.
Consequently convergence of all finite Gram words gives
\begin{equation}\label{eq:weighted-measure-convergence}
 (1-f_s)^2\,\zeta_L
       \longrightarrow (1-f(e^s))^2\,\mu
\end{equation}
against bounded continuous functions of finitely many coordinates in a
fixed window. The convergence is uniform when the window center stays
in a fixed-fraction array interior.

We need a finite-array bound for two dips before summing over windows.
Write \(\psi_L\) for the orthogonal derivative vector and
\(P_s=P_L(e^s)\), \(J_s=J_L(e^s)\). The differentiated eigen-equation is
\[
 (1-P_{s'})\psi_L=iJ_{s'}\Omega_L-a_{s'}\Omega_L
\]
for a scalar \(a_{s'}\). Since \(P_s\Omega_L=\Omega_L\),
\[
 (1-P_s)(1-P_{s'})\psi_L=-i[P_s,J_{s'}]\Omega_L.
\]
The spectral theorem therefore gives the exact identity
\[
 \int(1-f_s)^2(1-f_{s'})^2\,\dd\zeta_L
       =\|[P_s,J_{s'}]\Omega_L\|^2.
\]
Lemma~\ref{lem:row-normalization} bounds every singly marked row,
uniformly at all centers in a fixed-fraction array interior, once the
array is sufficiently long and its spacing sufficiently small.
In finite RTT the normalized intertwiner at separation \(s-s'\)
is an invertible diagonal matrix plus
\(O(e^{-\eta|s-s'|})\); the diagonal matrix commutes with the mark
\(I\otimes Z\). Expanding the conjugation error in auxiliary matrix
units expresses the commutator as a fixed finite sum of two-row products.
The single-row bounds consequently imply
\begin{equation}\label{eq:finite-two-dip}
 \int(1-f_s)^2(1-f_{s'})^2\,\dd\zeta_L
       \le C e^{-2\eta|s-s'|}.
\end{equation}
The constant is independent of the array length, small spacing and both
bulk centers. In particular there is no additive convergence error whose
sum over a growing number of pairs could accumulate.

\begin{lemma}[Packing angular windows]\label{lem:detector-packing}
The plane orbit mass satisfies
\begin{equation}\label{eq:kappa-upper}
 \kappa\le\frac{\eta}{2\lambda(\pi-\lambda)}.
\end{equation}
\end{lemma}
\begin{proof}
Choose finitely many real samples in a window, and let
\(Q(f)=\sum_j(1-f_{s_j})^2\).
Any bounded continuous detector \(F\) supported away from the all-ones
point can be written
\[
 F=\sum_j \frac{F}{Q}(1-f_{s_j})^2,
\]
with each quotient extended continuously by zero near that point.
Thus \eqref{eq:weighted-measure-convergence} gives convergence for such
detectors without assuming convergence of the full unweighted mass.
Similarly, products of two detectors are bounded by finite sums of
two-dip factors. The constants depend on the fixed window and mesh of
samples, not on the array center in the permitted bulk.

The cone condition and the dip estimate for inner functions provide
a fixed dip width and depth around each zero scale, independent of the
degree. Choose a fixed fine sampling mesh and a continuous detector
\(0\le F_T\le1\) in a window of length \(T\), equal to one when a
prescribed interior dip is detected and zero near all ones.
Each dilation orbit contributes an interval of centers of length at
least \(T-C\). Hence
\[
 \int F_T\,\dd\mu\ge(T-C)\kappa.
\]
Take the detector to be a continuous cutoff of
\(\max_j|1-f_{s_j}|\) at a fixed positive dip threshold \(a\).
Then \(F_T\le C_a\sum_j(1-f_{s_j})^2\), with \(C_a\) independent of
\(T\). Equation~\eqref{eq:finite-two-dip} bounds the product for two
windows separated by a gap \(g\) by \(CT^2e^{-2\eta g}\).
For \(N\) regularly spaced windows, their pair integrals sum to at most
\[
 \frac{C N T^2e^{-2\eta g}}{1-e^{-2\eta(T+g)}}.
\]

Pack windows of length \(T\) and gaps \(g\) into a bulk interval of
logarithmic length \((1-\gamma)B/\eta\).
Their number \(N_B\) satisfies
\[
 \frac{N_B}{B}\longrightarrow
          \frac{1-\gamma}{\eta(T+g)}.
\]
For numbers \(a_j\in[0,1]\),
\(\sum_j a_j\le1+\sum_{j<k}a_ja_k\).
Integrating this inequality against \(\zeta_L\), using bulk-uniform
detector convergence and the finite two-dip estimate, gives
\[
 \frac{1-\gamma}{\eta(T+g)}(T-C)\kappa
 \le \lim_{d\downarrow0}\limsup_{B\to\infty}\frac{I_L}{B}
       +\frac{CT^2 e^{-2\eta g}}{T+g}.
\]
For fixed \(T,g,\gamma\), all detectors use only a fixed finite number
of samples, so this application does not ask for convergence of a
growing word. The uniform estimates allow the inequalities to be summed
over their translated windows. The single-window convergence error,
after division by \(B\), is \((N_B/B)\varepsilon_{d,B,T}\), where
\(\lim_{d\downarrow0}\limsup_{B\to\infty}\varepsilon_{d,B,T}=0\).
Shortening the bulk interval at its ends by \(O(T+g)\) keeps every
sample in the permitted interior and does not change the limiting count.

Use Lemma~\ref{lem:overlap-asymptotic}, choose
\(g=A\log T\) with \(A\) sufficiently large, let \(T\to\infty\), and
then let \(\gamma\downarrow0\). This proves
\eqref{eq:kappa-upper}.
\end{proof}

\begin{proof}[Proof of Proposition~\ref{prop:normalization}]
Since \(\eta=\lambda/\pi\), the upper bound
\eqref{eq:kappa-upper} equals the lower bound
\eqref{eq:kappa-lower}. Thus \eqref{eq:kappa-value} holds.
The covariance conversion follows directly from \eqref{eq:green}.
\end{proof}

\subsection{Counting-coordinate accuracy and the finite trace limit}
\label{sec:finite-trace}

It remains to prove the overlap asymptotic used above. This subsection
and Section~\ref{sec:overlap-symbol} together prove
Lemma~\ref{lem:overlap-asymptotic}. We first compare the finite quadratic
logarithmic traces, divided by the physical length \(B\), with their bulk
convolution traces; the next subsection evaluates the resulting symbols.
All Taylor coefficients are taken at finite \(L\) before \(B\to\infty\)
at fixed \(d\), and then \(d\downarrow0\), with \(B=Ld/2\).

We first justify the precision needed to differentiate
\eqref{eq:exact-overlap}. The leading root density alone would not suffice.
In a dense interior interval, Lemmas~\ref{lem:root-limits}
and~\ref{lem:gaudin-inverse} give
\[
 \xi'(u)=d^{-1}+O(1),\qquad
 |\xi^{(m)}(u)|\le C_m d^{-1}\quad(m\ge1),
\]
with fixed derivative order. In the coordinate \(z=\xi(u)\), roots have
unit midpoint spacing and derivatives of the inverse coordinate satisfy
\[
 \left|\frac{\dd^m u}{\dd z^m}\right|\le C_m d^m.
\]
For a smooth exponentially localized test \(F\), localized away from the
ends, this gives
\begin{equation}\label{eq:counting-derivative}
 \int\left|
       \frac{\dd^m}{\dd z^m}F(u(z))\right|\,\dd z
       \le C_m d^{m-1}.
\end{equation}
Cutoffs can be put arbitrarily many decay lengths from the test center;
their errors tend to zero as \(B\to\infty\) at fixed \(d\).

Midpoint summation minus integration can be integrated by parts against
bounded periodic primitives of the unit-spaced counting discrepancy.
Using \eqref{eq:counting-derivative} to arbitrarily high fixed order gives
an error \(O_N(d^N)\) for every fixed \(N\), besides the vanishing
end errors. Apply this to the dressed density equation and all its
fixed derivatives. The uniform site array has the same quadrature
accuracy, so convolution inversion improves the density estimate to
\begin{equation}\label{eq:all-order-density}
 \xi'(u)=d^{-1}+O_N(d^N)
\end{equation}
with the same accuracy for its fixed derivatives in the interior.

Differentiating the Gaudin equation for the twist gives the constant
interior displacement
\[
 y_j(t)-y_j(0)
   = d z_*+O_N(d^N),\qquad
 z_*=\frac{t}{2(\pi-\lambda)},
\]
to each fixed Taylor order at zero. For the first derivative this is
\[
 \frac{d}{\pi(1+\widehat K(0))}
       =\frac{d}{2(\pi-\lambda)}.
\]
For higher derivatives, the differentiated site-array source is
arbitrarily small by midpoint quadrature, and difference terms vanish
on a constant displacement. Induction with the Gaudin inverse proves the
assertion. The statement concerns spacings and root differences, not an
arbitrarily chosen absolute mesh origin. Errors over several spacings
can grow polynomially with index separation; the kernels against which
they are summed decrease exponentially.

We next establish the finite matrix estimates needed for the trace
comparison, including the sparse ends. Write
\[
 x_j=y_j(0),\quad \rho(u)=\xi_0'(u),\quad
 \rho_j=\rho(x_j),\quad \ell_j=\rho_j^{-1},\quad
 \mathsf D_0=\diag(\ell_j),\quad \mathcal A=G_0\mathsf D_0=I+K\ell.
\]
For each matrix \(M(t)\) among \(M_{AC}(t),M_{CA}(t),G_C(t),G_A\),
use \(N(t)=M(t)\mathsf D_0\), keeping \(\mathsf D_0\) fixed at zero twist.
Then \(N(0)=\mathcal A\), and \(\log\det N-\log\det M\) is constant.
Equivalently, \(N(0)^{-1}N^{(r)}(0)
=\mathsf D_0^{-1}G_0^{-1}M^{(r)}(0)\mathsf D_0\), so the finite trace identity is
\begin{equation}\label{eq:second-log-trace}
 \left.\frac{\dd^2}{\dd t^2}\log\det M(t)\right|_{t=0}
 =\Tr\bigl(\mathcal A^{-1}N''(0)
       -\mathcal A^{-1}N'(0)\mathcal A^{-1}N'(0)\bigr).
\end{equation}
In particular, no derivative of \(\ell_j\) occurs, and an edge weight
has not been replaced by \(d\).

Here and below constants are uniform for sufficiently small \(d\) and
\(B\ge B_0(d)\), as in Section~\ref{sec:thermodynamic}.
Let \(p(u)=\sum_\alpha r(u-b_\alpha)\).
The sawtooth estimate \eqref{eq:counting-sawtooth} gives
\(\rho=p+O(1)\), with bounded integrated error.
The \(d/2\)-spaced array and exponential decay of \(r\) give
\(p(u)\le C/d\) on the line. Between the outer bare sites,
at least \(c/d\) sites lie within distance one of \(u\), once \(B\ge2\);
hence \(\rho(u)\ge c/d\) there for small \(d\).
At an exterior root, \eqref{eq:edge-index-density} instead gives
\(\rho_j\ge c\). Consecutive counting levels differ by one, so
\begin{equation}\label{eq:overlap-finite-array-bounds}
 \begin{gathered}
 0<\rho(u)\le C/d,\qquad cd\le\ell_j\le C,
 \qquad x_{j+1}-x_j\ge cd,\\
 \sup_u\sum_j\ell_j e^{-\gamma|u-x_j|}\le C_\gamma
 \quad(\gamma>0).
 \end{gathered}
\end{equation}
The last estimate is the local quadrature-mass bound of
Lemma~\ref{lem:root-limits}, summed over unit intervals.
Since \(K<0\), \(\rho\) is a sum of positive kernels \(s\) and
\(k=-K\); their fixed derivatives are bounded by constant multiples
of themselves. Thus \(|\rho^{(m)}(u)|\le C_m\rho(u)\).

Put \(v_j=y'_j(0)\), \(w_j=y''_j(0)\) on the twisted branch, and
\(\sigma(u)=\sum_\alpha s(u-b_\alpha)\).
Twice differentiating the finite counting equations gives
\[
 G_0v=\pi^{-1}\bbone,\qquad G_0w=-H,\qquad
 H_i=\sigma'(x_i)v_i^2
       -\sum_jK'(x_i-x_j)(v_i-v_j)^2.
\]
Since \(G_0^{-1}=\mathsf D_0\mathcal A^{-1}\),
Lemma~\ref{lem:gaudin-inverse} first gives \(|v_j|\le C\ell_j\).
Using \(\sigma(x_i)+\sum_j k(x_i-x_j)=\rho_i\), we then obtain
\[
 |H_i|\le C\left(\ell_i^2\rho_i+
                         \sum_j k(x_i-x_j)\ell_j^2\right)\le C.
\]
Here \(\rho_i\ell_i=1\), \(\ell_i\le C\), and the last bound in
\eqref{eq:overlap-finite-array-bounds} were used. Consequently
\begin{equation}\label{eq:overlap-root-jets}
 |v_j|+|w_j|\le C\ell_j.
\end{equation}
For the moving counting density \(\rho_t=\xi_{C(t)}'\), derivatives
at fixed \(u\) are
\[
 \dot\rho(u)=\sum_jK'(u-x_j)v_j,\qquad
 \ddot\rho(u)=\sum_j\{K'(u-x_j)w_j-K''(u-x_j)v_j^2\}.
\]
These functions and their required fixed spatial derivatives are
\(O(1)\), by \eqref{eq:overlap-finite-array-bounds} and
\eqref{eq:overlap-root-jets}. This includes all terminal indices.

We remove the coincident-column singularity before differentiation.
For either overlap orientation, let \(Y(t)\) be its on-shell list,
\(U(t)\) its column list, and
\(Q(t,u)=e^{-2it_Y}q_{Y(t)}(u)\), where \(t_Y=0\) for \(M_{AC}\)
and \(t_Y=t\) for \(M_{CA}\). Thus \(Q(t,Y_j(t))=-1\).
Set
\[
 h_j=U_j-Y_j,\qquad
 L_j(t)=\int_0^1\rho_{Y(t)}(Y_j+\theta h_j)\dd\theta,
 \qquad \mathcal E(z)=\frac{e^z-1}{z},\quad \mathcal E(0)=1.
\]
The identity \(q'/q=2\pi i\rho_Y\) gives
\(Q(t,U_j)=-e^{2\pi i h_jL_j}\). With
\[
 H_\lambda(u)=\frac{\coth u-\coth(u+i\lambda)}{-2\pi i},
 \qquad \beta_j(t)=1-e^{2\pi i h_jL_j},
\]
the finite scalar-product matrix is exactly
\(M_{ij}=K(u_{ij})+\beta_jH_\lambda(u_{ij})\),
where \(u_{ij}=U_j-Y_i\). Its constant coth tails have canceled.
On the diagonal this becomes the analytic identity
\begin{equation}\label{eq:overlap-coalescent-column}
 M_{jj}=K(h_j)+L_j\mathcal E(2\pi i h_jL_j)
       [h_j\coth h_j-h_j\coth(h_j+i\lambda)].
\end{equation}
At zero twist, \(h_j=0\), \(L_j=\rho_j\), and
\(|h'_j|+|h''_j|\le C\ell_j\).
Differentiating the integral defining \(L_j\), using the preceding
density bounds, gives \(|L'_j|+|L''_j|\le C\): every spatial
density derivative is multiplied by a velocity or acceleration
\(O(\ell_j)\), and the fixed-position time derivatives are \(O(1)\).
For \(z_j=2\pi i h_jL_j\), this implies
\[
 z'_j(0)=2\pi i\rho_jh'_j(0)=O(1),\qquad
 z''_j(0)=2\pi i(\rho_jh''_j+2h'_jL'_j)\big|_{t=0}=O(1).
\]
The bracket in \eqref{eq:overlap-coalescent-column} is analytic
with value one at zero; it and \(\mathcal E(z_j)\) have bounded first two
Taylor derivatives. The product rule and \(\ell_j\rho_j=1\) therefore
give \(\ell_j(|M'_{jj}(0)|+|M''_{jj}(0)|)\le C\).
No reciprocal displacement or reciprocal root velocity is introduced.

For \(i\ne j\), ordinary differentiation at \(u=x_j-x_i\) gives
\[
 \begin{split}
 M'_{ij}&=K'(u)u'_{ij}+\beta'_jH_\lambda(u),\\
 M''_{ij}&=K''(u)(u'_{ij})^2+K'(u)u''_{ij}
       +\beta''_jH_\lambda(u)+2\beta'_jH'_\lambda(u)u'_{ij}.
 \end{split}
\]
Here \(\beta_j(0)=0\), \(|\beta'_j(0)|+|\beta''_j(0)|\le C\),
and \(u'_{ij},u''_{ij}\) are bounded by \eqref{eq:overlap-root-jets}.
For some fixed \(\gamma>0\),
\[
 |K^{(r)}(u)|\le C_re^{-\gamma|u|},\quad
 |H_\lambda(u)|\le C(1+|u|^{-1})e^{-\gamma|u|},\quad
 |H'_\lambda(u)|\le C(1+|u|^{-2})e^{-\gamma|u|}.
\]
Choose small \(\alpha>0\) and write
\(\|T\|_{\alpha,x}=\sup_i\sum_j
 e^{\alpha|x_i-x_j|}|T_{ij}|\).
The minimum gap bounds the two pole factors by \(C/d\) and \(C/d^2\).
Multiplication by the column weight \(\ell_j\), local quadrature mass,
and the diagonal estimate prove
\begin{equation}\label{eq:overlap-taylor-row-bounds}
 \|M'(0)\mathsf D_0\|_{\alpha,x}\le C/d,\qquad
 \|M''(0)\mathsf D_0\|_{\alpha,x}\le C/d^2
 \quad(M=M_{AC},M_{CA}).
\end{equation}
For the Gaudin denominator, put \(r_i(t)=\rho_t(y_i(t))\).
The formulas
\[
 r'_i=\dot\rho_i+\rho'_iv_i,\qquad
 r''_i=\ddot\rho_i+2\dot\rho'_iv_i+\rho''_iv_i^2+\rho'_iw_i
\]
give bounded first two derivatives. Differentiating
\(G_{ij}(t)=r_i(t)\delta_{ij}+K(y_i(t)-y_j(t))\) and summing with
the same column weights thus gives
\(\|G'_C(0)\mathsf D_0\|_{\alpha,x}+\|G''_C(0)\mathsf D_0\|_{\alpha,x}\le C\).
Lemma~\ref{lem:gaudin-inverse}, with both signs of its exponential
weight, gives \(\|\mathcal A^{-1}\|_{\alpha,x}\le C\).
Weighted absolute row norms are submultiplicative by the triangle
inequality. Hence the matrix inside the trace in
\eqref{eq:second-log-trace}, and each of its diagonal entries,
are bounded by \(C/d^2\), uniformly in \(B\).

The discarded count includes all exterior roots. Indeed, with
\(T(u)=\int_u^\infty\rho(v)\dd v\), the half-integer tail levels
and integrated sawtooth estimate give
\[
 \#\{j:x_j\ge x\}\le T(x)+1
 \le\sum_\alpha\bar r(x-b_\alpha)+C,
 \qquad \bar r(v)=\int_v^\infty r(u)\dd u.
\]
For \(x=B/2-R\), at most \(C(R/d+1)\) bare sites lie to its right,
each contributing at most \(1/2\). The other sites contribute at most
\(C\sum_{m\ge0}e^{-md/(2\eta)}\le C/d\).
Reflection therefore proves
\begin{equation}\label{eq:overlap-end-count}
 \#\{j:|x_j|\ge B/2-R\}\le C(R+1)/d.
\end{equation}
Deleting these diagonal entries costs at most
\(C(R+1)/(d^3B)\) after division by \(B\).
This tends to zero at fixed \(d\) if \(R=R(B)\to\infty\) and
\(R(B)/B\to0\); a fixed number of nested strips has the same property.

For the inverse comparison we use one common weighted row algebra,
including at the sparse ends:
\begin{equation}\label{eq:overlap-common-row-norm}
 \|T\|_{a,d}=\sup_i\sum_j e^{ad|i-j|}|T_{ij}|,
 \qquad T_n=I+(dK(d(i-j)))_{1\le i,j\le n}.
\end{equation}
Since \(|x_i-x_j|\ge cd|i-j|\), the actual-matrix bounds above hold
in this norm for sufficiently small fixed \(a>0\).
Choose \(a\) also so that \(\int e^{a|u|}|K(u)|\dd u<1\).
For small \(d\), \(\|T_n-I\|_{a,d}\le q<1\); the finite Neumann
series gives \(\|T_n^{-1}\|_{a,d}\le(1-q)^{-1}\), uniformly in \(n\).
Thus the exact identity
\[
 \mathcal A^{-1}-T_n^{-1}
       =\mathcal A^{-1}(T_n-\mathcal A)T_n^{-1}
\]
holds in the same weighted algebra for both inverses.
Two nested physical cutoffs separated by \(R\) inside the bare array
have at least \(cR/d-O(1)\) intervening root indices, since
\(\rho\ge c/d\) there. Any index path from a retained inner diagonal
to the outer strip therefore gains a factor \(Ce^{-c'R}\).

In the interior, \eqref{eq:all-order-density} and the twist-jet
estimates give \(\ell_j=d+O_N(d^N)\) and arbitrarily accurate
comparisons of root differences with \(d(j-i)\).
Apply these to the normalized Taylor matrices, using
\eqref{eq:overlap-coalescent-column} on the diagonal.
Fixed derivatives at distinct-root poles cost only powers of \(1/d\);
polynomial growth with index separation is absorbed by exponential
decay. Choosing a higher initial accuracy order therefore makes the
interior row errors of the Taylor matrices and of
\(\mathcal A-T_n\) \(O_N(d^N)\) for every fixed \(N\), besides errors
vanishing at fixed \(d\) as
the strips recede. Their global row bounds have only polynomial
costs in \(1/d\). The inverse identity and telescoping the fixed
products in \eqref{eq:second-log-trace}, with nested strips, give a
retained diagonal error at most
\[
 P(d^{-1})\bigl[O_N(d^N)+o_{B\to\infty,d}(1)+C_de^{-c'R(B)}\bigr]
\]
for a fixed polynomial \(P\). No intermediate row sum costs a factor
growing with \(B\). Since exactly \(n=B/d\) indices occur, the trace
per physical length costs only one additional factor \(1/d\).
Send \(B\to\infty\) first, and choose \(N>\deg P+1\) before
letting \(d\downarrow0\).

Finally, the weighted Neumann series compares \(T_n^{-1}\) with the
full-lattice convolution inverse: every omitted path from an interior
row crosses the matrix boundary, and so has the same exponentially
small bound. The same argument applies to the finite products.
The corresponding convolution trace integrand has
diagonal coefficient \(d(2\pi)^{-1}\int_{-\pi/d}^{\pi/d}
\widehat F(p)\dd p\). Multiplication by \(n/B=1/d\) gives precisely
the Fourier measure \(\dd p/(2\pi)\) per unit physical length.
No infinite-volume trace-class assertion is involved.

\subsection{Evaluation of the symbol}
\label{sec:overlap-symbol}
\begin{proof}[Proof of Lemma~\ref{lem:overlap-asymptotic}]
In the product part of \eqref{eq:exact-overlap}, the quadratic coefficient
is a sum of
\((\log[x+i\lambda])''\) times the two first displacements. Its integral
on the real line is \(2\), from the two limits of \(\coth(x+i\lambda)\).
The constant accelerations cancel between cross and within-set terms.
With root density \(1/d\) and displacement \(dz_*\), the coefficient
per unit length is \(2z_*^2\), with a vanishing two-limit error.

For the determinants, the bulk sequence of \(dM_{AC}\), at row-column
difference \(y=d(k+z_*)\), is \(d\) times
\begin{equation}\label{eq:overlap-toeplitz}
 \frac{(1-e^{2\pi i z_*})\coth y-\coth(y-i\lambda)
                         +e^{2\pi i z_*}\coth(y+i\lambda)}
      {-2\pi i}.
\end{equation}
For \(M_{CA}\) replace \(z_*\) by \(-z_*\); for the two Gaudin
denominators use \(z_*=0\). The coincident pole is continued by its
removable limit. The preceding trace comparison reduces the computation
to the quadratic coefficient of these bulk symbols.

Use Fourier exponent \(-ipd(k+z_*)\).
The simple pole contributes the shifted-sampling identity
\[
 \sum_{k\in\Z}\frac{e^{-ipd(k+z_*)}}{k+z_*}
    =\pi\bigl(\cot(\pi z_*)-i\sgn p\bigr),
 \qquad 0<|p|<\pi/d,
\]
with symmetric summation. Subtract \(1/y\) from the singular cotangent
before treating the remaining part. Its derivative is integrable and
analytic in every strip
\(|\Im y|<a<\min(\lambda,\pi-\lambda)\).
Contour shift therefore makes every nonzero alias exponentially small
in \(1/d\), including its fixed Taylor coefficients. This remains
uniform at the band ends since a nonzero alias has absolute frequency
at least \(\pi/d\); the value at \(p=0\) is irrelevant to the integral.

Write \(a_0=e^{-\lambda p}\) and \(b_0=e^{-(\pi-\lambda)p}\).
The continuous transform together with the sampled pole gives
\[
 \frac{(1-e^{2\pi i z_*})(1+a_0b_0)/2-b_0
                              +e^{2\pi i z_*}a_0}
      {1-a_0b_0}
       +\frac{1+e^{2\pi i z_*}}2
 =\frac{(1-b_0)(1+e^{2\pi i z_*}a_0)}
        {1-a_0b_0}.
 \]
Dividing by the value at zero twist, the relative symbol is
\[
 \frac{1+e^{2\pi i z_*}a_0}{1+a_0}.
\]
The offset Fourier phases cancel between the two orientations.
Their combined quadratic logarithmic coefficient is
\[
 -4\pi^2 z_*^2\frac{a_0}{(1+a_0)^2}.
\]
Since
\[
 \int_{\R}\frac{e^{-\lambda p}}{(1+e^{-\lambda p})^2}\,
       \frac{\dd p}{2\pi}=\frac1{2\pi\lambda},
\]
the determinant contribution per length is \(-2\pi z_*^2/\lambda\).
For each finite \(L\), let \([t^2]\) denote the Taylor coefficient at
\(t=0\) of the local logarithm fixed above. Adding the product contribution
proves the explicit coefficient limit
\[
 \lim_{d\downarrow0}\limsup_{\substack{B\to\infty\\ B=Ld/2}}
 \left|\frac1B[t^2]\log|\langle\Omega_{L,0},\Omega_{L,t}\rangle|^2
       +\frac1{2\lambda(\pi-\lambda)}\right|=0.
\]
The Taylor coefficient is taken before either array limit; no common
analytic neighborhood or uniform finite-\(t\) expansion is asserted.
By \eqref{eq:overlap-quadratic} that coefficient is \(-I_L\), which
proves the lemma.
\end{proof}

